\section{Implementation}
\label{sec:implementation}

Both implementations are written in Go (1.24+) and share the same
\rfcjcs-compliant JSON canonicalization pipeline.  They differ only in
the digit-generation module:

\begin{itemize}
\item \texttt{impl-schubfach}: uses the \schubfach algorithm with
      \texttt{bits.Mul64} for 128-bit multiplication.  Zero heap
      allocation in the formatting hot path.

\item \texttt{impl-json-canon}: uses the \burgerdybvig algorithm with
      \texttt{math/big.Int}.  Uses \texttt{sync.Pool} to reuse
      big-integer buffers.
\end{itemize}

Both implementations:
\begin{itemize}
\item Do not use \texttt{encoding/json} for canonicalization.
\item Produce static binaries (\texttt{CGO\_ENABLED=0}).
\item Pass identical conformance gates (154 structural cases + 286{,}362
      oracle vectors).
\item Implement identical CLI ABI: \texttt{canonicalize} and
      \texttt{verify} modes with defined exit codes and error classes.
\end{itemize}

The shared pipeline handles JSON parsing, recursive key sorting (UTF-16
code unit order per \rfcjcs), string escaping, and output serialization.
The digit-generation module is invoked only for JSON number values.
